\documentclass{icontec}
\usepackage[utf8]{inputenc}
\usepackage[T1]{fontenc}
\usepackage{geometry}
\usepackage{graphicx}
\usepackage{float}
\usepackage{hyperref}
\usepackage[english,indonesian]{babel}

\title{Optimasi Pemrosesan Template LaTeX Menggunakan Pendekatan Berbasis AI}
\author{Academic Support Agent}
\date{\today}

\begin{document}
\maketitle

\begin{abstract}

\end{abstract}

\chapter{PENDAHULUAN}


\chapter{TINJAUAN PUSTAKA}


% Original content (commented out):
% \documentclass{icontec}

\usepackage{blindtext}
\blindmathtrue

\title{TÍTULO DEL TRABAJO O PROYECTO}
\author{NOMBRES Y APELLIDOS COMPLETOS DEL AUTOR O AUTORES}
\legend{LEYENDA DEL TRABAJO\\
(clase del trabajo realizado, tesis, monografía)}
\director{DIRECTOR, PRESIDENTE, ASESOR, PROFESOR\\
Nombres y apellidos completos de la persona que dirigió el trabajo con su titulo académico o cargo}
\institution{NOMBRE DE LA INSTITUCIÓN}
\faculty{FACULTAD, DEPARTAMENTO, DEPENDENCIA, SECCIÓN O ÁREA}
\city{CIUDAD}
\date{AÑO}

\begin{document}

\maketitle

\makeapprobation

\tableofcontents

\listoffigures

\listoftables

\Blinddocument

\end{document}


\end{document}